\chapter{Pile technique et versions}

\section{Technologies principales}
L'application repose sur une pile \textbf{React/Vite/TypeScript} moderne avec
\textbf{Tailwind CSS} pour le style et \textbf{Supabase} pour le backend.

\begin{longtable}{p{4.2cm}p{2.2cm}p{8cm}}
\toprule
\textbf{Technologie} & \textbf{Version} & \textbf{Rôle} \\
\midrule
\endhead
React & 18.3 & Bibliothèque d'interface (composants, hooks). \\
Vite & 5.4 & Bundler / serveur de développement rapide. \\
TypeScript & 5.8 & Typage statique de tout le code. \\
Tailwind CSS & 3.4 & Système de design utilitaire (tokens sémantiques). \\
shadcn/ui (Radix) & --- & Composants accessibles (dialogues, menus, onglets\dots). \\
React Router DOM & 6.30 & Routage côté client (SPA). \\
TanStack React Query & 5.83 & Cache et synchronisation des requêtes serveur. \\
Supabase JS & 2.97 & Client base de données / Auth / Storage / Functions. \\
\bottomrule
\end{longtable}

\section{Bibliothèques notables}
\begin{longtable}{p{4.5cm}p{10cm}}
\toprule
\textbf{Paquet} & \textbf{Usage dans le projet} \\
\midrule
\endhead
\texttt{i18next} / \texttt{react-i18next} & Internationalisation FR/AR + détection de langue. \\
\texttt{katex}, \texttt{react-katex}, \texttt{rehype-katex}, \texttt{remark-math} & Rendu des formules mathématiques (\LaTeX) dans les leçons et le chat. \\
\texttt{react-markdown}, \texttt{remark-gfm}, \texttt{rehype-raw} & Rendu du contenu Markdown (cours, solutions, commentaires IA). \\
\texttt{@tiptap/*} & Éditeur de texte riche WYSIWYG pour l'espace éditorial. \\
\texttt{recharts} & Graphiques de progression (tableaux de bord, analytics). \\
\texttt{framer-motion} & Animations (compagnon-bot, transitions). \\
\texttt{jspdf}, \texttt{html2canvas} & Export PDF (cours, rapports parents). \\
\texttt{react-hook-form} + \texttt{zod} & Formulaires et validation de schémas. \\
\texttt{@dnd-kit/*} & Glisser-déposer (réordonnancement de sections éditoriales). \\
\texttt{date-fns} & Manipulation des dates. \\
\texttt{lucide-react}, \texttt{react-icons} & Jeux d'icônes. \\
\texttt{sonner}, \texttt{vaul}, \texttt{embla-carousel} & Toasts, tiroirs, carrousels. \\
\bottomrule
\end{longtable}

\section{Outils de développement et qualité}
\begin{itemize}
  \item \textbf{ESLint 9} + \texttt{typescript-eslint} --- qualité et cohérence du code.
  \item \textbf{Vitest 3} + \texttt{@testing-library/react} --- tests unitaires/composants.
  \item \textbf{Playwright 1.57} --- tests de bout en bout (navigateur).
  \item \textbf{PostCSS / Autoprefixer} --- chaîne CSS.
  \item \textbf{@vitejs/plugin-react-swc} --- compilation rapide via SWC.
\end{itemize}

\section{Scripts npm}
\begin{lstlisting}[language=bash,caption={package.json --- scripts}]
npm run dev        # serveur de developpement Vite
npm run build      # build de production
npm run build:dev  # build en mode developpement
npm run lint       # analyse ESLint
npm run test       # tests Vitest (une passe)
npm run test:watch # tests en mode surveillance
\end{lstlisting}

\begin{notebox}
Le projet est construit avec l'écosystème \textbf{Lovable}. Les fonctions Edge
sont déployées automatiquement à chaque modification ; le fichier \filepath{.env}
est auto-rempli avec \texttt{VITE\_SUPABASE\_URL},
\texttt{VITE\_SUPABASE\_PUBLISHABLE\_KEY} et \texttt{VITE\_SUPABASE\_PROJECT\_ID}.
\end{notebox}
